\documentclass[12pt,letterpaper]{article}

\usepackage[utf8]{inputenc}
\usepackage[T1]{fontenc}
\usepackage[spanish,es-tabla]{babel}
\usepackage{lmodern}
\usepackage{amssymb}
\usepackage[margin=2.4cm]{geometry}
\usepackage{array}
\usepackage{longtable}
\usepackage[table]{xcolor}
\usepackage{graphicx}
\usepackage{float}
\usepackage{caption}
\usepackage{enumitem}
\usepackage{fancyhdr}
\usepackage{tikz}
\usetikzlibrary{arrows.meta,positioning,shapes.geometric,babel}
\usepackage[hidelinks]{hyperref}

\graphicspath{{capturas/}}

\definecolor{azul}{HTML}{0F2D57}
\definecolor{azulMedio}{HTML}{1F5FBF}
\definecolor{azulClaro}{HTML}{E8F0FC}
\renewcommand{\arraystretch}{1.32}
\setlength{\parindent}{0pt}
\setlength{\parskip}{0.55em}
\captionsetup{font=small,labelfont={bf,color=azul}}

\setlength{\headheight}{15pt}
\pagestyle{fancy}
\fancyhf{}
\fancyhead[L]{\small\color{azul} Sistema web de inventario MIN}
\fancyhead[R]{\small\color{azul} Desarrollo de Requisitos de Software}
\fancyfoot[C]{\thepage}

% Captura de pantalla: \captura[ancho]{archivo}{pie}{etiqueta}
\newcommand{\captura}[4][0.95]{%
  \begin{figure}[H]
    \centering
    \fbox{\includegraphics[width=#1\linewidth,height=0.72\textheight,keepaspectratio]{#2}}
    \caption{#3}
    \label{#4}
  \end{figure}}

% Encabezado de tablas de requisitos
\newcommand{\encabezado}[3]{\rowcolor{azulClaro} \textbf{#1} & \textbf{#2} & \textbf{#3} \\ \hline}

\begin{document}

% ======================================================= PORTADA
\begin{titlepage}
  \centering
  {\large\bfseries Universidad Autónoma de Ciudad Juárez\par}
  {\normalsize Instituto de Ingeniería y Tecnología\par}
  \vspace{0.4cm}
  {\color{azul}\rule{\linewidth}{1.2pt}}\par
  \vspace{2.2cm}
  {\Large Desarrollo de Requisitos de Software\par}
  \vspace{1.2cm}
  {\color{azul}\Huge\bfseries Sistema web de control de inventario\par}
  \vspace{0.4cm}
  {\Large Caso de estudio: Mantenimiento Industrial del Norte (MIN)\par}
  \vspace{2.4cm}
  \begin{tabular}{rl}
    \textbf{Alumno:} & Alejandro Martínez Contreras \\
    \textbf{Matrícula:} & \underline{\hspace{4cm}} \\
    \textbf{Profesor(a):} & \underline{\hspace{4cm}} \\
    \textbf{Modalidad:} & Individual \\
  \end{tabular}
  \vfill
  {\color{azul}\rule{\linewidth}{1.2pt}}\par
  {\large Ciudad Juárez, Chihuahua \quad·\quad Noviembre de 2026\par}
\end{titlepage}

\tableofcontents
\clearpage

% ======================================================= INTRODUCCIÓN
\section{Introducción}

Una empresa dedicada al mantenimiento de maquinaria industrial controla manualmente las piezas y refacciones de su almacén. Esta práctica provoca diferencias entre las existencias reales y las registradas, pérdida de información y dificultad para saber qué técnico utilizó cada material y con qué autorización.

En este documento se especifican los requisitos de un sistema web que resuelve ese problema y se presenta la aplicación desarrollada para cumplirlos. El sistema, llamado \textbf{MIN – Control de Inventario}, organiza el trabajo en tres roles: el \textbf{Administrador} mantiene el catálogo de piezas y las cuentas de usuario; el \textbf{Gerente} consulta el inventario, autoriza o rechaza solicitudes y genera reportes; y el \textbf{Técnico} consulta piezas disponibles, solicita materiales y da seguimiento a sus solicitudes.

Cada salida de material queda ligada a una solicitud, a un técnico y a un gerente que la aprobó. Las existencias se actualizan automáticamente y todo cambio se registra como movimiento de entrada, salida o ajuste. Además, el sistema muestra alertas de piezas agotadas o con existencias por debajo del mínimo.

\subsection{Objetivo}

Especificar, desarrollar y validar una aplicación web que permita administrar el inventario de piezas industriales con control de acceso por rol, flujo de solicitud y autorización de materiales, actualización automática de existencias, alertas de stock y trazabilidad de movimientos.

\subsection{Alcance}

El sistema cubre el almacén de la empresa: catálogo de piezas, entradas por recepción de proveedor, salidas autorizadas mediante solicitudes, ajustes de existencias, alertas, reportes y administración de usuarios. Quedan fuera de este alcance las compras a proveedores, la facturación y la integración con otros sistemas de la empresa.

\subsection{Tecnologías utilizadas}

\begin{itemize}[leftmargin=1.4em]
  \item \textbf{HTML5 y plantillas Jinja2:} estructura de las vistas.
  \item \textbf{CSS3:} diseño visual adaptable a computadora y teléfono, sin bibliotecas externas.
  \item \textbf{JavaScript:} confirmaciones antes de eliminar e impresión de reportes.
  \item \textbf{Python 3 y Flask 3:} servidor web, rutas, reglas de negocio y control de acceso.
  \item \textbf{SQLite:} base de datos relacional que conserva piezas, usuarios, solicitudes y movimientos.
  \item \textbf{Git y GitHub:} control de versiones y publicación del código fuente.
  \item \textbf{Playwright:} recorrido automatizado de la aplicación en Google Chrome para obtener las capturas de evidencia.
\end{itemize}

% ======================================================= ACTORES
\section{Actores y permisos}

La Tabla~\ref{tab:permisos} resume qué puede hacer cada rol. Los permisos se validan en el servidor: si un usuario intenta abrir una función que no le corresponde, el sistema responde con un error 403.

\begin{table}[H]
\centering
\small
\begin{tabular}{|p{6.8cm}|c|c|c|}
\hline
\rowcolor{azulClaro} \textbf{Función} & \textbf{Administrador} & \textbf{Gerente} & \textbf{Técnico} \\ \hline
Iniciar y cerrar sesión & \checkmark & \checkmark & \checkmark \\ \hline
Ver dashboard y alertas de inventario & \checkmark & \checkmark & \checkmark \\ \hline
Consultar y buscar piezas & \checkmark & \checkmark & \checkmark \\ \hline
Registrar, modificar y eliminar piezas & \checkmark & & \\ \hline
Registrar entradas de inventario & \checkmark & & \\ \hline
Administrar usuarios & \checkmark & & \\ \hline
Crear solicitudes de material & & & \checkmark \\ \hline
Revisar el estado de sus solicitudes & & & \checkmark \\ \hline
Aprobar o rechazar solicitudes & & \checkmark & \\ \hline
Consultar movimientos de inventario & \checkmark & \checkmark & \\ \hline
Generar reportes y exportar CSV & \checkmark & \checkmark & \\ \hline
\end{tabular}
\caption{Matriz de permisos por rol.}
\label{tab:permisos}
\end{table}

% ======================================================= REQUISITOS FUNCIONALES
\section{Requisitos funcionales}

Los requisitos funcionales describen los servicios que el sistema debe proporcionar a cada tipo de usuario. Se redactan con el verbo «deber» e incluyen un criterio de verificación que permite comprobarlos con las evidencias de la Sección~\ref{sec:evidencias}.

\begin{longtable}{|>{\bfseries}p{1.3cm}|p{8.3cm}|p{5.2cm}|}
\hline
\encabezado{ID}{Requisito}{Criterio de verificación}
\endfirsthead
\hline
\encabezado{ID}{Requisito}{Criterio de verificación}
\endhead
RF01 & El sistema debe permitir iniciar y cerrar sesión a usuarios con rol de Administrador, Gerente o Técnico, cada uno con usuario y contraseña propios. & Las tres cuentas entran al sistema; una contraseña incorrecta muestra un mensaje de error. \\ \hline
RF02 & El sistema debe mostrar a cada rol únicamente las opciones de menú y acciones que le corresponden, y bloquear el acceso directo a funciones no autorizadas. & Un técnico que abre \texttt{/usuarios} recibe el error 403. \\ \hline
RF03 & El sistema debe mostrar un dashboard con el total de piezas y unidades, el número de piezas en stock, con stock bajo y sin stock, una gráfica de distribución, las alertas, las solicitudes relevantes y los últimos movimientos. & El dashboard presenta los indicadores actualizados para cada rol. \\ \hline
RF04 & El sistema debe contar con al menos 20 piezas industriales registradas con código, nombre, categoría, cantidad, stock mínimo y ubicación. & El catálogo inicial contiene 24 piezas con todos sus datos. \\ \hline
RF05 & El sistema debe permitir consultar las piezas y buscarlas por código, nombre o ubicación, así como filtrarlas por categoría y por estado de existencias. & La búsqueda «sensor» devuelve únicamente los sensores. \\ \hline
RF06 & El sistema debe permitir al Administrador registrar, modificar y eliminar piezas, validando código único, campos obligatorios y cantidades no negativas. & La pieza nueva aparece en el catálogo; un código repetido se rechaza. \\ \hline
RF07 & El sistema debe permitir al Administrador registrar entradas de inventario indicando pieza, cantidad y referencia (factura o proveedor). & La existencia aumenta y se registra un movimiento de entrada. \\ \hline
RF08 & El sistema debe permitir al Técnico crear solicitudes de material indicando pieza, cantidad y motivo u orden de trabajo, sin exceder la existencia disponible. & La solicitud recibe un folio y queda en estado «En revisión». \\ \hline
RF09 & El sistema debe permitir al Técnico consultar el estado de sus solicitudes (en revisión, aprobada o rechazada) y el comentario del gerente. & El técnico ve sus solicitudes con el estado y la resolución. \\ \hline
RF10 & El sistema debe permitir al Gerente aprobar o rechazar solicitudes pendientes y registrar un comentario. & La solicitud cambia de estado y guarda quién y cuándo la resolvió. \\ \hline
RF11 & Al aprobar una solicitud, el sistema debe descontar automáticamente la cantidad autorizada de las existencias; si no hay existencias suficientes, la aprobación debe rechazarse. & La existencia de la pieza disminuye exactamente en la cantidad aprobada. \\ \hline
RF12 & El sistema debe mostrar alertas de piezas agotadas (cantidad igual a 0) y con existencias bajas (cantidad menor o igual al mínimo). & Las piezas en alerta aparecen en el dashboard, en la campana y en los filtros. \\ \hline
RF13 & El sistema debe registrar cada movimiento de entrada, salida o ajuste con fecha, pieza, cantidad, existencia resultante, usuario y referencia. & La aprobación de una solicitud genera una salida con su folio. \\ \hline
RF14 & El sistema debe permitir al Gerente y al Administrador generar reportes del inventario y exportar el inventario y los movimientos a archivos CSV. & El reporte muestra totales y se descargan los archivos CSV. \\ \hline
RF15 & El sistema debe permitir al Administrador crear, modificar, desactivar y eliminar cuentas de usuario y asignarles un rol. & La nueva cuenta puede iniciar sesión con el rol asignado. \\ \hline
\end{longtable}

% ======================================================= REQUISITOS NO FUNCIONALES
\section{Requisitos no funcionales}

Los requisitos no funcionales establecen condiciones de seguridad, integridad, usabilidad, rendimiento y portabilidad que el sistema debe cumplir.

\begin{longtable}{|>{\bfseries}p{1.5cm}|p{2.4cm}|p{6.2cm}|p{4.7cm}|}
\hline
\rowcolor{azulClaro} \textbf{ID} & \textbf{Atributo} & \textbf{Requisito} & \textbf{Criterio de verificación} \\ \hline
\endfirsthead
\hline
\rowcolor{azulClaro} \textbf{ID} & \textbf{Atributo} & \textbf{Requisito} & \textbf{Criterio de verificación} \\ \hline
\endhead
RNF01 & Seguridad & Las contraseñas deben almacenarse como hash y nunca en texto claro. & La tabla de usuarios solo contiene valores hash. \\ \hline
RNF02 & Seguridad & Todas las páginas, excepto el inicio de sesión, deben requerir una sesión activa, y los permisos de rol deben validarse en el servidor. & El acceso sin sesión redirige al login; el acceso con un rol no autorizado devuelve 403. \\ \hline
RNF03 & Seguridad & Los formularios que modifican datos deben incluir un token CSRF y las consultas a la base de datos deben ser parametrizadas. & Una petición sin token válido se rechaza con error 400. \\ \hline
RNF04 & Integridad & El descuento de existencias debe ser atómico: una solicitud no puede aprobarse dos veces ni dejar existencias negativas. & Una segunda aprobación de la misma solicitud se rechaza y la existencia no cambia. \\ \hline
RNF05 & Persistencia & La información debe conservarse en una base de datos SQLite después de cerrar la aplicación. & Los registros siguen disponibles al reiniciar el servidor. \\ \hline
RNF06 & Trazabilidad & Ningún cambio de existencias debe ocurrir sin un movimiento registrado. & La suma de movimientos de cada pieza coincide con su existencia. \\ \hline
RNF07 & Usabilidad & La interfaz debe ser clara, en español, con colores que distingan el estado del stock (verde, amarillo, rojo) y mensajes de confirmación o error tras cada acción. & Cada acción muestra un mensaje y los estados se identifican por color y texto. \\ \hline
RNF08 & Adaptabilidad & La interfaz debe poder utilizarse en pantallas de escritorio y de teléfonos móviles. & Las vistas se reorganizan sin perder información en un ancho de 390 px. \\ \hline
RNF09 & Rendimiento & Las consultas y búsquedas deben responder en menos de 2 segundos con el volumen de datos de un almacén mediano. & Las pantallas cargan de inmediato en el equipo local. \\ \hline
RNF10 & Portabilidad & La aplicación debe ejecutarse con Python 3.10 o posterior, Flask 3.x y un navegador moderno, sin servicios externos. & El sistema se instala con \texttt{pip install -r requirements.txt}. \\ \hline
RNF11 & Mantenibilidad & El código debe estar versionado en GitHub, documentado en un README y contar con una prueba automatizada del flujo principal. & \texttt{python smoke\_test.py} termina sin errores. \\ \hline
\end{longtable}

% ======================================================= HISTORIAS
\section{Historias de usuario}

\begin{longtable}{|>{\bfseries}p{1.3cm}|p{8.1cm}|p{5.4cm}|}
\hline
\encabezado{ID}{Historia de usuario}{Criterio de aceptación}
\endfirsthead
\hline
\encabezado{ID}{Historia de usuario}{Criterio de aceptación}
\endhead
HU01 & Como \textbf{administrador}, quiero registrar y actualizar las piezas del almacén para mantener un catálogo confiable. & La pieza guardada aparece con código, nombre, categoría, cantidad y ubicación. \\ \hline
HU02 & Como \textbf{administrador}, quiero registrar la recepción de material para que las existencias reflejen lo que llega del proveedor. & La existencia aumenta y queda un movimiento de entrada con la referencia. \\ \hline
HU03 & Como \textbf{administrador}, quiero crear cuentas y asignar roles para controlar quién usa el sistema. & La nueva cuenta inicia sesión y solo ve las funciones de su rol. \\ \hline
HU04 & Como \textbf{técnico}, quiero buscar una pieza y ver cuántas unidades hay para saber si puedo usarla en una reparación. & La búsqueda muestra la existencia y el estado de la pieza. \\ \hline
HU05 & Como \textbf{técnico}, quiero solicitar material indicando la orden de trabajo para que el gerente lo autorice. & La solicitud recibe un folio y queda «En revisión». \\ \hline
HU06 & Como \textbf{técnico}, quiero revisar el estado de mis solicitudes para saber si ya puedo recoger el material. & El técnico ve si su solicitud fue aprobada o rechazada y el comentario. \\ \hline
HU07 & Como \textbf{gerente}, quiero aprobar o rechazar solicitudes para controlar el consumo de material. & Al aprobar, la existencia disminuye; al rechazar, no cambia. \\ \hline
HU08 & Como \textbf{gerente}, quiero ver alertas y reportes para decidir qué piezas reabastecer. & El reporte lista las piezas en alerta y la cantidad sugerida a comprar. \\ \hline
\end{longtable}

% ======================================================= ANÁLISIS Y DISEÑO
\section{Análisis y diseño}

\subsection{Diagrama de casos de uso}

\begin{figure}[H]
\centering
\resizebox{0.95\linewidth}{!}{%
\begin{tikzpicture}[
    actor/.style={draw=azul,fill=azulClaro,rounded corners=3pt,minimum width=2.6cm,minimum height=1cm,align=center,font=\bfseries},
    usecase/.style={draw=azul,ellipse,fill=white,minimum width=4.2cm,minimum height=0.85cm,align=center,font=\small},
    line/.style={draw=azul,thick}]
  \draw[azul,rounded corners=6pt,thick] (2.6,1.2) rectangle (11.4,-11.2);
  \node[font=\bfseries,text=azul] at (7,0.8) {Sistema MIN};
  \node[actor] (admin) at (0,-1.6) {Administrador};
  \node[actor] (ger) at (0,-6.4) {Gerente};
  \node[actor] (tec) at (14,-5) {Técnico};
  \node[usecase] (login) at (7,0) {Iniciar / cerrar sesión};
  \node[usecase] (dash) at (7,-1) {Consultar dashboard y alertas};
  \node[usecase] (cons) at (7,-2) {Consultar y buscar piezas};
  \node[usecase] (piezas) at (7,-3) {Administrar piezas};
  \node[usecase] (entrada) at (7,-4) {Registrar entrada};
  \node[usecase] (users) at (7,-5) {Administrar usuarios};
  \node[usecase] (aut) at (7,-6.2) {Autorizar / rechazar solicitud};
  \node[usecase] (mov) at (7,-7.3) {Consultar movimientos};
  \node[usecase] (rep) at (7,-8.3) {Generar reportes};
  \node[usecase] (sol) at (7,-9.4) {Crear solicitud};
  \node[usecase] (est) at (7,-10.4) {Revisar estado de solicitud};
  \foreach \caso in {login,dash,cons,piezas,entrada,users,mov,rep} \draw[line] (admin.east) -- (\caso.west);
  \foreach \caso in {login,dash,cons,aut,mov,rep} \draw[line] (ger.east) -- (\caso.west);
  \foreach \caso in {login,dash,cons,sol,est} \draw[line] (tec.west) -- (\caso.east);
\end{tikzpicture}}
\caption{Diagrama de casos de uso.}
\label{fig:casos-uso}
\end{figure}

\subsection{Diagrama de actividades: solicitud de piezas}

El proceso central del sistema inicia con la solicitud del técnico y termina con la actualización del inventario y el registro del movimiento.

\begin{figure}[H]
\centering
\resizebox{0.9\linewidth}{!}{%
\begin{tikzpicture}[
  flow/.style={-Latex,thick,draw=azul},
  step/.style={draw=azul,rounded corners=3pt,fill=azulClaro,minimum width=3.6cm,minimum height=0.85cm,align=center,font=\small},
  decision/.style={draw=azul,diamond,aspect=2.2,fill=white,inner sep=1pt,align=center,font=\small},
  terminal/.style={draw=azul,ellipse,fill=azulClaro,minimum width=2cm,minimum height=0.7cm,align=center,font=\small}]
  \node[terminal] (inicio) at (0,0) {Inicio};
  \node[step] (crear) at (0,-1.3) {Técnico crea solicitud\\(pieza, cantidad, motivo)};
  \node[decision] (disp) at (0,-3.1) {¿Cantidad $\le$\\existencia?};
  \node[step] (err1) at (5,-3.1) {Mostrar error\\al técnico};
  \node[step] (pend) at (0,-4.9) {Solicitud «En revisión»};
  \node[step] (rev) at (0,-6.3) {Gerente revisa la solicitud};
  \node[decision] (dec) at (0,-8.1) {¿Aprueba?};
  \node[step] (rech) at (-5,-8.1) {Estado «Rechazada»\\(sin cambio de stock)};
  \node[decision] (stock) at (0,-10) {¿Stock\\suficiente?};
  \node[step] (err2) at (5,-10) {Avisar existencias\\insuficientes};
  \node[step] (desc) at (0,-11.8) {Descontar existencias\\y marcar «Aprobada»};
  \node[step] (mov) at (0,-13.2) {Registrar movimiento\\de salida con folio};
  \node[decision] (alerta) at (0,-15) {¿Cantidad $\le$\\mínimo?};
  \node[step] (al) at (5,-15) {Mostrar alerta\\de stock};
  \node[terminal] (fin) at (0,-16.8) {Fin};
  \draw[flow] (inicio) -- (crear);
  \draw[flow] (crear) -- (disp);
  \draw[flow] (disp) -- node[above,font=\scriptsize] {No} (err1);
  \draw[flow] (err1.north) |- (crear.east);
  \draw[flow] (disp) -- node[right,font=\scriptsize] {Sí} (pend);
  \draw[flow] (pend) -- (rev);
  \draw[flow] (rev) -- (dec);
  \draw[flow] (dec) -- node[above,font=\scriptsize] {No} (rech);
  \draw[flow] (rech.south) |- (fin.west);
  \draw[flow] (dec) -- node[right,font=\scriptsize] {Sí} (stock);
  \draw[flow] (stock) -- node[above,font=\scriptsize] {No} (err2);
  \draw[flow] (stock) -- node[right,font=\scriptsize] {Sí} (desc);
  \draw[flow] (desc) -- (mov);
  \draw[flow] (mov) -- (alerta);
  \draw[flow] (alerta) -- node[above,font=\scriptsize] {Sí} (al);
  \draw[flow] (al.south) |- (fin.east);
  \draw[flow] (alerta) -- node[right,font=\scriptsize] {No} (fin);
\end{tikzpicture}}
\caption{Diagrama de actividades del proceso de solicitud, autorización y actualización.}
\label{fig:actividades}
\end{figure}

\subsection{Modelo de datos}

La base de datos tiene cuatro tablas. Cada solicitud se relaciona con el técnico que la crea, el gerente que la resuelve y la pieza solicitada. Cada movimiento conserva el código y nombre de la pieza, el usuario responsable y la existencia resultante.

\begin{figure}[H]
\centering
\resizebox{\linewidth}{!}{%
\begin{tikzpicture}[
  class/.style={draw=azul,thick,fill=azulClaro,rounded corners=2pt,inner sep=7pt,align=left,font=\small},
  rel/.style={thick,draw=azul}]
  \node[class] (usuario) at (0,0) {\begin{tabular}{l}
    \textbf{Usuario}\\ \hline
    id: Integer\\ usuario: String\\ nombre: String\\
    rol: Administrador $|$ Gerente $|$ Técnico\\
    password\_hash: String\\ activo: Boolean
  \end{tabular}};
  \node[class] (pieza) at (12,0) {\begin{tabular}{l}
    \textbf{Pieza}\\ \hline
    id: Integer\\ codigo: String (único)\\ nombre: String\\ categoria: String\\
    cantidad: Integer $\ge 0$\\ minimo: Integer\\ ubicacion: String\\ actualizado: DateTime
  \end{tabular}};
  \node[class] (sol) at (6,-6.2) {\begin{tabular}{l}
    \textbf{Solicitud}\\ \hline
    id: Integer (folio SOL-\#\#\#)\\ cantidad: Integer $>0$\\ motivo: String\\
    estado: Pendiente $|$ Aprobada $|$ Rechazada\\
    creada: DateTime\\ resuelta: DateTime\\ comentario: String
  \end{tabular}};
  \node[class] (mov) at (17,-6.2) {\begin{tabular}{l}
    \textbf{Movimiento}\\ \hline
    id: Integer\\ tipo: Entrada $|$ Salida $|$ Ajuste\\ cantidad: Integer\\
    existencia: Integer\\ referencia: String\\ fecha: DateTime
  \end{tabular}};
  \draw[rel] (usuario.south) -- node[left,font=\scriptsize,align=center] {1 técnico\\0..*} (sol.north west);
  \draw[rel,dashed] (usuario.south east) -- node[above,font=\scriptsize,sloped] {0..1 gerente resuelve} (sol.north);
  \draw[rel] (pieza.south west) -- node[right,font=\scriptsize] {1 — 0..*} (sol.north east);
  \draw[rel] (pieza.south east) -- node[right,font=\scriptsize] {1 — 0..*} (mov.north);
  \draw[rel] (usuario.east) -- node[above,font=\scriptsize] {1 registra 0..* movimientos} (pieza.west);
\end{tikzpicture}}
\caption{Diagrama de clases del modelo de datos.}
\label{fig:modelo}
\end{figure}

% ======================================================= EVIDENCIAS
\clearpage
\section{Evidencias de funcionamiento}
\label{sec:evidencias}

Las capturas se obtuvieron ejecutando la aplicación en un servidor local y recorriéndola en Google Chrome con cada uno de los tres usuarios. Todas las cifras que aparecen corresponden a operaciones realizadas en el propio sistema durante la prueba.

\subsection{Inicio de sesión}

\captura[0.8]{01_login.png}{Pantalla de inicio de sesión común para los tres roles (RF01).}{fig:login}
\captura[0.8]{02_login_error.png}{Credenciales incorrectas: el sistema no permite el acceso y muestra un error (RF01).}{fig:login-error}

\subsection{Funciones del Administrador}

\captura[0.85]{03_admin_dashboard.png}{Dashboard del Administrador: accesos rápidos, resumen general, gráfica de existencias, alertas, solicitudes por autorizar y últimos movimientos (RF03, RF12).}{fig:admin-dash}
\captura{04_admin_piezas.png}{Catálogo con las 24 piezas industriales: código, nombre, categoría, cantidad, mínimo, ubicación, estado y acciones de edición, entrada y eliminación (RF04, RF06).}{fig:admin-piezas}
\captura{05_admin_agregar_pieza.png}{Registro de una nueva pieza (RF06).}{fig:admin-nueva}
\captura{06_admin_pieza_registrada.png}{Confirmación del alta: la pieza P025 aparece en el catálogo (RF06).}{fig:admin-registrada}
\captura{07_admin_editar_pieza.png}{Edición de una pieza existente; un cambio de cantidad genera un movimiento de ajuste (RF06, RF13).}{fig:admin-editar}
\captura{08_admin_recibir_inventario.png}{Recepción de inventario: entrada de 10 motores DC con referencia de factura (RF07).}{fig:admin-entrada}
\captura{09_admin_entrada_registrada.png}{La entrada queda registrada en la bitácora de movimientos con la existencia resultante (RF07, RF13).}{fig:admin-entrada-ok}
\captura{10_admin_usuarios.png}{Administración de usuarios y roles (RF15).}{fig:admin-usuarios}
\captura{11_admin_usuario_creado.png}{Alta de un nuevo técnico por parte del Administrador (RF15).}{fig:admin-usuario-nuevo}

\subsection{Funciones del Técnico}

\captura[0.85]{12_tecnico_dashboard.png}{Dashboard del Técnico: el menú solo muestra Dashboard, Piezas y Solicitudes (RF02, RF03).}{fig:tec-dash}
\captura{13_tecnico_busqueda.png}{Búsqueda de piezas por texto («sensor»); el técnico solo puede consultar y solicitar (RF05).}{fig:tec-busqueda}
\captura{16_tecnico_acceso_denegado.png}{Un técnico que intenta abrir la administración de usuarios recibe el error 403 (RF02, RNF02).}{fig:tec-403}

\subsection{Funciones del Gerente}

\captura[0.85]{17_gerente_dashboard.png}{Dashboard del Gerente con las solicitudes por autorizar (RF03, RF10).}{fig:ger-dash}
\captura[0.85]{24_gerente_reportes.png}{Reporte de inventario: entradas, salidas, solicitudes, existencias por categoría, piezas con más salidas, solicitudes por técnico y piezas a reabastecer, con exportación a CSV (RF14).}{fig:ger-reportes}

% ======================================================= PROCESO COMPLETO
\section{Evidencia del proceso de solicitud, autorización y actualización}

Esta sección sigue una misma solicitud de principio a fin: el técnico Jorge Ramírez solicita 5 sensores inductivos M18 (código P003) para la orden de trabajo OT-2311.

\textbf{Paso 1. El técnico crea la solicitud.}
\captura{14_tecnico_nueva_solicitud.png}{Formulario de solicitud: pieza, cantidad y motivo (RF08).}{fig:p1}
\captura{15_tecnico_solicitud_enviada.png}{La solicitud SOL-004 queda «En revisión»; aún no se descuenta el inventario (RF08, RF09).}{fig:p2}

\textbf{Paso 2. Existencia antes de la autorización.}
\captura{18_stock_antes.png}{Antes de la aprobación, la pieza P003 tiene 15 unidades.}{fig:p3}

\textbf{Paso 3. El gerente revisa y autoriza.}
\captura{19_gerente_por_autorizar.png}{El gerente ve las solicitudes pendientes, la existencia disponible y escribe un comentario (RF10).}{fig:p4}
\captura{20_gerente_solicitud_aprobada.png}{SOL-004 aprobada: el sistema confirma el descuento de 5 unidades. También se observa una solicitud rechazada (SOL-003) que no modificó las existencias (RF10, RF11).}{fig:p5}

\textbf{Paso 4. El sistema actualiza el inventario automáticamente.}
\captura{21_stock_despues.png}{Después de la aprobación, la pieza P003 tiene 10 unidades ($15 - 5$) (RF11).}{fig:p6}

\textbf{Paso 5. Se registra el movimiento.}
\captura{22_movimiento_salida.png}{Movimiento de salida ligado al folio SOL-004, con el usuario que autorizó y la existencia resultante (RF13).}{fig:p7}
\captura[0.85]{23_movimientos.png}{Bitácora completa de entradas, salidas y ajustes (RF13).}{fig:p8}

\textbf{Paso 6. El técnico consulta el resultado.}
\captura{27_tecnico_estado_solicitudes.png}{El técnico ve su solicitud aprobada, quién la autorizó y el comentario (RF09).}{fig:p9}

\subsection{Alertas de inventario}

\captura{25_alertas_stock_bajo.png}{Filtro de piezas con existencias bajas (RF12).}{fig:alerta-bajo}
\captura{26_alertas_agotadas.png}{Filtro de piezas agotadas (RF12).}{fig:alerta-sin}

\subsection{Vista en dispositivo móvil}

\captura[0.4]{28_vista_movil.png}{Dashboard del Gerente en un teléfono de 390 px de ancho (RNF08).}{fig:movil}

% ======================================================= TRAZABILIDAD
\section{Matriz de trazabilidad}

\begin{longtable}{|>{\bfseries}p{1.4cm}|p{7.2cm}|p{6.2cm}|}
\hline
\encabezado{Req.}{Descripción breve}{Evidencia}
\endfirsthead
\hline
\encabezado{Req.}{Descripción breve}{Evidencia}
\endhead
RF01 & Inicio de sesión de tres roles & Figuras \ref{fig:login}, \ref{fig:login-error}, \ref{fig:admin-dash}, \ref{fig:tec-dash}, \ref{fig:ger-dash} \\ \hline
RF02 & Permisos por rol & Figuras \ref{fig:tec-dash}, \ref{fig:tec-403} \\ \hline
RF03 & Dashboard & Figuras \ref{fig:admin-dash}, \ref{fig:tec-dash}, \ref{fig:ger-dash} \\ \hline
RF04 & 20 o más piezas registradas & Figura \ref{fig:admin-piezas} \\ \hline
RF05 & Consulta y búsqueda & Figura \ref{fig:tec-busqueda} \\ \hline
RF06 & Alta, edición y eliminación de piezas & Figuras \ref{fig:admin-piezas}--\ref{fig:admin-editar} \\ \hline
RF07 & Entradas de inventario & Figuras \ref{fig:admin-entrada}, \ref{fig:admin-entrada-ok} \\ \hline
RF08 & Solicitud del técnico & Figuras \ref{fig:p1}, \ref{fig:p2} \\ \hline
RF09 & Estado de solicitudes & Figuras \ref{fig:p2}, \ref{fig:p9} \\ \hline
RF10 & Aprobar o rechazar & Figuras \ref{fig:p4}, \ref{fig:p5} \\ \hline
RF11 & Actualización de existencias & Figuras \ref{fig:p3}, \ref{fig:p6} \\ \hline
RF12 & Alertas de stock & Figuras \ref{fig:admin-dash}, \ref{fig:alerta-bajo}, \ref{fig:alerta-sin} \\ \hline
RF13 & Registro de movimientos & Figuras \ref{fig:admin-entrada-ok}, \ref{fig:p7}, \ref{fig:p8} \\ \hline
RF14 & Reportes y CSV & Figura \ref{fig:ger-reportes} \\ \hline
RF15 & Administración de usuarios & Figuras \ref{fig:admin-usuarios}, \ref{fig:admin-usuario-nuevo} \\ \hline
\end{longtable}

Además, el archivo \texttt{smoke\_test.py} del repositorio verifica automáticamente el inicio de sesión, los permisos, el ciclo de vida de las piezas, la solicitud, la aprobación, el rechazo, el descuento de existencias, la protección CSRF, la imposibilidad de aprobar dos veces la misma solicitud y la exportación de reportes.

% ======================================================= REPOSITORIO
\section{Repositorio de GitHub}

El código fuente completo, las instrucciones de instalación, la prueba automatizada y las capturas de evidencia están disponibles en:

\begin{center}
\fcolorbox{azul}{azulClaro}{\parbox{0.85\linewidth}{\centering\large
\url{https://github.com/TU-USUARIO/Inventario_MIN}}}
\end{center}

Cuentas de demostración: \texttt{admin}, \texttt{gerente}, \texttt{tecnico1} y \texttt{tecnico2}; sus contraseñas iniciales están en el archivo \texttt{README.md}.

% ======================================================= CONCLUSIONES
\section{Conclusiones}

\subsection*{Conclusión individual — Alejandro Martínez Contreras}

Al desarrollar esta actividad comprobé que escribir primero los requisitos facilita mucho la construcción del sistema. El caso de estudio pedía cosas como «actualizar las existencias», pero al especificarlas tuve que decidir cuándo ocurre ese descuento (solo al aprobar la solicitud), qué pasa si no hay existencias suficientes y cómo evitar que una solicitud se apruebe dos veces. Esas decisiones se convirtieron en requisitos verificables y después en pruebas.

También aprendí que separar a los usuarios por roles no es solo esconder botones: los permisos deben validarse en el servidor, porque un usuario puede escribir la dirección de una página directamente. Por eso incluí la evidencia del error 403. La bitácora de movimientos resolvió el problema de no saber quién utiliza los materiales, ya que cada salida queda ligada a un folio, a un técnico y al gerente que la autorizó.

Finalmente, la matriz de trazabilidad me permitió comprobar que cada requisito del caso de estudio tiene una evidencia. Considero que este enfoque reduce los errores que la empresa tenía con el registro manual y deja una base clara para futuras mejoras, como órdenes de compra automáticas o lectura de códigos QR.

% ======================================================= REFERENCIAS
\begin{thebibliography}{99}

\bibitem{sommerville}
Sommerville, I. (2016). \textit{Software Engineering} (10.ª ed.). Pearson.

\bibitem{pressman}
Pressman, R. S. y Maxim, B. R. (2020). \textit{Software Engineering: A Practitioner's Approach} (9.ª ed.). McGraw-Hill Education.

\bibitem{wiegers}
Wiegers, K. y Beatty, J. (2013). \textit{Software Requirements} (3.ª ed.). Microsoft Press.

\bibitem{iso29148}
ISO/IEC/IEEE. (2018). \textit{ISO/IEC/IEEE 29148:2018 Systems and software engineering — Life cycle processes — Requirements engineering}. International Organization for Standardization.

\bibitem{iso25010}
ISO/IEC. (2023). \textit{ISO/IEC 25010:2023 Systems and software engineering — Systems and software Quality Requirements and Evaluation (SQuaRE) — Product quality model}. International Organization for Standardization.

\bibitem{cohn}
Cohn, M. (2004). \textit{User Stories Applied: For Agile Software Development}. Addison-Wesley.

\bibitem{flask}
Pallets Projects. (2025). \textit{Flask Documentation (3.1.x)}. \url{https://flask.palletsprojects.com/}

\bibitem{sqlite}
SQLite Consortium. (2025). \textit{SQLite Documentation}. \url{https://www.sqlite.org/docs.html}

\bibitem{mdn}
Mozilla. (2025). \textit{MDN Web Docs: HTML, CSS y JavaScript}. \url{https://developer.mozilla.org/es/}

\bibitem{owasp}
OWASP Foundation. (2024). \textit{Cross-Site Request Forgery Prevention Cheat Sheet}. \url{https://cheatsheetseries.owasp.org/cheatsheets/Cross-Site_Request_Forgery_Prevention_Cheat_Sheet.html}

\end{thebibliography}

\end{document}
